\section{Experimental Protocol}
\label{sec:parallel_experimental_protocol}
\hspace{\parindent}
The experimental protocol compares the sequential random-ordering baseline with the proposed
parallel framework. Both configurations were evaluated on the same three large-scale scenarios to
ensure that the observed differences result from the planning strategy rather than from changes in
the problem instances.

The sequential baseline evaluates random robot priority orderings with TEA* until either a complete
valid solution is found or the execution limit \(T_{\max}=15\,\mathrm{s}\) is reached. This value was
chosen as a practical upper bound, since even a \(15s\) planning interval is already
long for an online multi-robot coordination system. For each scenario, this process is repeated
over 100 independent runs. In each successful run, the time required to
obtain the first valid solution and its corresponding
cost are recorded. Runs that fail to find a complete valid
solution within \(T_{\max}\) are counted separately and excluded from the mean time and cost
calculations, so the reported averages are computed only over successful runs.

The proposed framework was evaluated using \(1\), \(3\), \(5\), and \(6\) worker threads. For each
run, both the initial ordering of heuristic groups in the priority queue and the ordering of
heuristics within each group were randomized, preventing the evaluation from manually favouring any
specific heuristic family for a given scenario. After this initialization step, the dispatching
mechanism follows the group usage counters described in Chapter~\ref{chap:ch7}, prioritizing
heuristics from less used groups throughout the planning cycle.

Each parallel configuration was also executed 100 times per scenario, with each run recording the
time to the first valid solution, the cost of the first valid solution, the time at which the best
solution was found, the cost of the best solution, and the time required to complete the evaluation
of the heuristic pool. In this context, the best-known solution corresponds to the lowest-cost
solution identified for the same scenario in the heuristic evaluation of Chapter~\ref{chap:ch6}.
Therefore, the time to the best solution measures how quickly the parallel framework reaches the
best solution previously obtained by the individual heuristic analysis.

In addition to the mean values reported in the tables, Empirical Cumulative Distribution Functions
(ECDFs) are used to complement the analysis of planning times across the 100 runs. These curves are
reported for the time to the first valid solution and for the time to the best-known solution,
allowing the consistency and convergence behaviour of each configuration to be compared beyond the
average values alone. The sequential baseline is included in these plots only when it provides a
meaningful number of valid runs.





\section{Parallel Framework Results}
\label{sec:parallel_framework_results}

\subsection{Scenario 1}
\label{subsec:parallel_results_scene1}
\hspace{\parindent}
The first large-scale scenario, described in
Section~\ref{subsubsec:large_scale_docking_to_warehouses}, involves movements from docking
areas towards warehouse regions. As shown in Table~\ref{tab:parallel_results_scene1}, the
sequential random-ordering baseline was unable to obtain a valid solution in any of the 100 runs.
This result is consistent with the heuristic evaluation reported for the same scenario in
Section~\ref{subsubsec:large_scale_docking_to_warehouses}, where random priority orderings were
shown to be ineffective for this problem instance. In contrast, all heuristic-based configurations
solved every run, showing that the developed framework substantially improves planning
robustness in this instance.


\begin{table}[H]
\centering
\caption{Parallel framework results for large-scale Scenario~1 over 100 runs.}
\label{tab:parallel_results_scene1}
\scriptsize
\setlength{\tabcolsep}{4pt}
\begin{tabular}{@{}lccccc@{}}
\toprule
\textbf{Configuration} &
\makecell{\textbf{Valid}\\\textbf{runs}} &
\makecell{\textbf{Mean time to}\\\textbf{first valid (s)}} &
\makecell{\textbf{Mean first}\\\textbf{cost (s)}} &
\makecell{\textbf{Mean time to}\\\textbf{best solution (s)}} &
\makecell{\textbf{Mean time to}\\\textbf{complete pool (s)}} \\
\midrule
Sequential random TEA* & 0/100   & --   & --       & --   & --    \\
Heuristics-based TEA*, 1 thread     & 100/100 & 2.88 & 48122.63 & 5.19 & 10.56 \\
Heuristics-based TEA*, 3 threads    & 100/100 & 1.93 & 48115.73 & 2.97 & 5.34  \\
Heuristics-based TEA*, 5 threads    & 100/100 & 1.74 & 48095.02 & 2.49 & 4.51  \\
Heuristics-based TEA*, 6 threads    & 100/100 & 1.91 & 48095.02 & 2.71 & 3.85  \\
\bottomrule
\end{tabular}
\end{table}

The timing distributions are presented in Figure~\ref{fig:scenario1_time_ecdfs}. Since the
sequential baseline did not produce any valid solution, it is omitted from the ECDF of the first
valid solution. The first-valid-solution ECDF shows that the single-threaded configuration can
produce very early solutions in some runs, but its curve rises more gradually, indicating lower
consistency across executions. The three-thread configuration shifts the distribution towards shorter
times, while the five-thread configuration reaches a large fraction of successful runs earlier than
the remaining configurations. The six-thread configuration follows a similar trend to the
five-thread case, but does not provide a visible advantage in this scenario. For the best-known
solution, the ECDFs show the same general behaviour, with the multi-threaded configurations
reaching high cumulative fractions substantially earlier than the single-threaded configuration.

\begin{figure}[H]
\centering
\begin{subfigure}{0.49\textwidth}
    \centering
    \includegraphics[width=\textwidth]{scenario_1_first_valid_time_ecdf.png}
    \caption{Time to first valid solution.}
    \label{fig:scenario1_first_valid_time_ecdf}
\end{subfigure}
\hfill
\begin{subfigure}{0.49\textwidth}
    \centering
    \includegraphics[width=\textwidth]{scenario_1_best_plan_time_ecdf.png}
    \caption{Time to best-known solution.}
    \label{fig:scenario1_best_plan_time_ecdf}
\end{subfigure}
\caption{Empirical cumulative distributions of planning times in Scenario~1 over 100 runs.}
\label{fig:scenario1_time_ecdfs}
\end{figure}

\subsection{Scenario 2}
\label{subsec:parallel_results_scene2}
\hspace{\parindent}
The second large-scale scenario, introduced in
Section~\ref{subsubsec:large_scale_movements_between_warehouses}, represents movements between
warehouse regions. Unlike Scenario~1, the sequential random-ordering baseline solved all 100 runs,
which allows a direct comparison between the baseline and the proposed framework. This behaviour
is consistent with the heuristic evaluation of the same scenario, where random priority orderings
were shown to be more effective than in the other large-scale instances. Even so, the results in
Table~\ref{tab:parallel_results_scene2} show that the heuristic-based configurations improve the
quality of the first valid solution, while the multi-threaded configurations also reduce the time
required to reach the best-known solution.

\begin{table}[H]
\centering
\caption{Parallel framework results for large-scale Scenario~2 over 100 runs.}
\label{tab:parallel_results_scene2}
\scriptsize
\setlength{\tabcolsep}{3.5pt}
\begin{tabular}{@{}lcccccc@{}}
\toprule
\textbf{Configuration} &
\makecell{\textbf{Valid}\\\textbf{runs}} &
\makecell{\textbf{Mean time to}\\\textbf{first valid (s)}} &
\makecell{\textbf{Mean first}\\\textbf{cost (s)}} &
\makecell{\textbf{Solution-time}\\\textbf{speedup}} &
\makecell{\textbf{Mean time to}\\\textbf{best solution (s)}} &
\makecell{\textbf{Mean time to}\\\textbf{complete pool (s)}} \\
\midrule
Sequential random TEA*          & 100/100 & 1.92 & 33296.28 & 1.00 & --    & --    \\
Heuristic-based TEA*, 1 thread  & 100/100 & 2.32 & 32855.42 & 0.83 & 10.52 & 15.17 \\
Heuristic-based TEA*, 3 threads & 100/100 & 1.38 & 32070.76 & 1.39 & 5.01  & 7.82  \\
Heuristic-based TEA*, 5 threads & 100/100 & 1.54 & 31435.03 & 1.25 & 4.26  & 6.41  \\
Heuristic-based TEA*, 6 threads & 100/100 & 1.70 & 31547.22 & 1.13 & 3.88  & 6.07  \\
\bottomrule
\end{tabular}
\end{table}

The timing distributions in Figure~\ref{fig:scenario2_time_ecdfs} show a different behaviour for the
first valid solution and for the best-known solution. In the first-valid-solution ECDF, the
three-thread configuration rises earlier than the remaining configurations, indicating that it reaches
a larger fraction of successful runs within shorter time intervals. The five- and six-thread
configurations remain competitive, but their curves are slightly shifted to the right, suggesting that
the additional parallelism does not translate directly into a faster first response in this scenario.
For the best-known solution, the curves become progressively more favourable as the number of
threads increases, showing that higher thread counts are more effective when the objective is to
reach the best solution found during the heuristic pool evaluation rather than only the first valid
one.

\begin{figure}[H]
\centering
\begin{subfigure}{0.49\textwidth}
    \centering
    \includegraphics[width=\textwidth]{scenario_2_first_valid_time_ecdf.png}
    \caption{Time to first valid solution.}
    \label{fig:scenario2_first_valid_time_ecdf}
\end{subfigure}
\hfill
\begin{subfigure}{0.49\textwidth}
    \centering
    \includegraphics[width=\textwidth]{scenario_2_best_plan_time_ecdf.png}
    \caption{Time to best-known solution.}
    \label{fig:scenario2_best_plan_time_ecdf}
\end{subfigure}
\caption{Empirical cumulative distributions of planning times in Scenario~2 over 100 runs.}
\label{fig:scenario2_time_ecdfs}
\end{figure}


\subsection{Scenario 3}
\label{subsec:parallel_results_scene3}
\hspace{\parindent}
The third large-scale scenario, described in
Section~\ref{subsubsec:large_scale_warehouses_to_docking}, considers movements from warehouse
regions back to docking areas. This scenario was also shown to be highly restrictive in the
heuristic evaluation, where random priority orderings rarely produced complete valid plans. The
same behaviour is observed in Table~\ref{tab:parallel_results_scene3}, where the sequential
random-ordering baseline obtained valid solutions in only two of the 100 runs. The corresponding
baseline averages are therefore computed only over these two successful executions and should be
interpreted together with the low success rate. In contrast, every heuristic-based configuration
solved all runs, confirming the robustness gain provided by the developed framework.

\begin{table}[H]
\centering
\caption{Parallel framework results for large-scale Scenario~3 over 100 runs.}
\label{tab:parallel_results_scene3}
\scriptsize
\setlength{\tabcolsep}{3.5pt}
\begin{tabular}{@{}lcccccc@{}}
\toprule
\textbf{Configuration} &
\makecell{\textbf{Valid}\\\textbf{runs}} &
\makecell{\textbf{Mean time to}\\\textbf{first valid (s)}} &
\makecell{\textbf{Mean first}\\\textbf{cost (s)}} &
\makecell{\textbf{Solution-time}\\\textbf{speedup}} &
\makecell{\textbf{Mean time to}\\\textbf{best solution (s)}} &
\makecell{\textbf{Mean time to}\\\textbf{complete pool (s)}} \\
\midrule
Sequential random TEA*          & 2/100   & 15.22 & 57313.03 & 1.00 & --    & --    \\
Heuristic-based TEA*, 1 thread  & 100/100 & 7.32  & 52818.71 & 2.08 & 15.42 & 34.32 \\
Heuristic-based TEA*, 3 threads & 100/100 & 3.41  & 52324.14 & 4.46 & 6.09  & 16.03 \\
Heuristic-based TEA*, 5 threads & 100/100 & 2.83  & 49858.53 & 5.37 & 4.15  & 12.41 \\
Heuristic-based TEA*, 6 threads & 100/100 & 2.89  & 47778.68 & 5.27 & 3.82  & 11.83 \\
\bottomrule
\end{tabular}
\end{table}

The ECDFs in Figure~\ref{fig:scenario3_time_ecdfs} emphasize the difference between random
priority sampling and heuristic-guided parallel exploration. The sequential random baseline is
omitted from the first-valid-solution plot because only two valid runs were obtained. Among the
heuristic-based configurations, the single-threaded curve rises much more slowly, indicating that
evaluating one ordering at a time leads to a less consistent response. The three-thread
configuration already shifts the distribution substantially to the left, while the five- and six-thread
configurations concentrate most runs within shorter time intervals. For the best-known solution,
the separation between configurations becomes clearer, with higher thread counts reaching large
cumulative fractions earlier and therefore reducing the time required to converge to the best-known
solution.

\begin{figure}[H]
\centering
\begin{subfigure}{0.49\textwidth}
    \centering
    \includegraphics[width=\textwidth]{scenario_3_first_valid_time_ecdf.png}
    \caption{Time to first valid solution.}
    \label{fig:scenario3_first_valid_time_ecdf}
\end{subfigure}
\hfill
\begin{subfigure}{0.49\textwidth}
    \centering
    \includegraphics[width=\textwidth]{scenario_3_best_plan_time_ecdf.png}
    \caption{Time to best-known solution.}
    \label{fig:scenario3_best_plan_time_ecdf}
\end{subfigure}
\caption{Empirical cumulative distributions of planning times in large-scale Scenario~3 over 100 runs. The sequential random baseline is omitted from the first-valid-solution plot because only two valid solutions were obtained.}
\label{fig:scenario3_time_ecdfs}
\end{figure}

\section{Overall Performance Discussion}
\label{sec:parallel_overall_performance}
